\section{Related Work} \label{sec:7related-work}

\newcommand{\cmk}{\textcolor{xhs_green}{\ding{51}}}
\newcommand{\xmk}{\textcolor{xhs_red}{\ding{55}}}

\begin{table}[t]
\centering
\caption{Comparison with prior symbolic and probabilistic approaches to protocols with compromisable keys. See legend below for the meaning of criteria C1--C5.}
\label{tab:related-work-comparison}
\small
\renewcommand{\arraystretch}{1.2}
\begin{tabularx}{\columnwidth}{@{}>{\raggedright\arraybackslash}X ccccc@{}}
\toprule
\textbf{Work} & \textbf{C1} & \textbf{C2} & \textbf{C3} & \textbf{C4} & \textbf{C5} \\
\midrule
Shi et al.~\cite{shi2023blesc}                       & \xmk & \xmk & \xmk & \cmk & \xmk \\
Claverie et al.~\cite{claverie2023tamarin}           & \xmk & \xmk & \xmk & \cmk & \xmk \\
Jangid et al.~\cite{jangid2023bluepe}                & \xmk & \xmk & \xmk & \cmk & \xmk \\
Cheval et al.~\cite{cheval2023hash}                  & \cmk & \xmk & \xmk & \xmk & \cmk \\
Cremers et al.~\cite{cremers2019prime}               & \cmk & \xmk & \xmk & \xmk & \cmk \\
Siedlecka-Lamch et al.~\cite{siedlecka2016probabilistic} & \cmk & \xmk & \xmk & --- & \xmk \\
\midrule
\textbf{EntropyVerif (ours)}                         & \cmk & \cmk & \cmk & \cmk & \cmk \\
\bottomrule
\end{tabularx}

\vspace{0.4em}
\begin{flushleft}
\footnotesize
\textbf{C1}~Protocol-agnostic modeling. \hspace{0.4em}
\textbf{C2}~Compound-key decomposition. \hspace{0.4em}
\textbf{C3}~Graded attacker with bit-level bound. \hspace{0.4em}
\textbf{C4}~No kernel modification (\cmk{} only if the approach does not recompile \Tamarin{} or alter the term-rewriting engine; ``---'' marks Siedlecka-Lamch et al., who use PRISM instead). \hspace{0.4em}
\textbf{C5}~New protocol-level attack discovered.
\end{flushleft}
\end{table}

Table~\ref{tab:related-work-comparison} positions our work against five lines of prior research. We next elaborate on each.

\noindent\textbf{Empirical studies of low-entropy vulnerabilities.}
Early studies~\cite{gong1993protecting,kelsey1997secure} highlighted the profound impact of low-entropy keys on cryptographic-protocol security.
In practice, a long list of severe vulnerabilities has been attributed to low-entropy cryptographic variables, including those disclosed in~\cite{jangid2023bluepe,schneier1998cryptanalysis,ferguson1999cryptographic,gpgcracking} in addition to the KNOB and PMKID attacks examined in this paper.
These studies identify the problem but do not offer an analysis methodology; our work provides the first automated tool that discovers such vulnerabilities at the protocol design stage.

\noindent\textbf{Enhanced attacker models for specific primitives.}
Several works extend symbolic tools with enhanced attackers against specific cryptographic primitives, such as weak hash functions~\cite{cheval2023hash}, unsafe Diffie-Hellman groups~\cite{cremers2019prime}, malleable digital signatures~\cite{jackson2019sign}, and AEAD schemes~\cite{cremers2023aead}.
These works target primitive-level weaknesses rather than key-length-induced weaknesses, and typically require modifications to the tool's internal term-rewriting logic.
Brute-force attacks across arbitrary compound keys are not covered by any of them and remain, to the best of our knowledge, unstudied in symbolic models prior to this work.

\noindent\textbf{Protocol-specific reveal rules.}
Shi et al.~\cite{shi2023blesc} proposed an ad-hoc approach to discover brute-force-associated attacks in BLE, where a fixed reveal rule allows the attacker to compromise one particular key.
Claverie et al.~\cite{claverie2023tamarin} and Jangid et al.~\cite{jangid2023bluepe} follow a similar protocol-specific pattern to enable compromise of a single cryptographic variable in BLE pairing.
These approaches are tailored to a known attack and cannot generalize, because (i) the reveal rule hard-codes one specific function symbol, (ii) there is no notion of compound-key decomposition, and (iii) no graded-attacker model constrains which keys are actually within reach of a realistic adversary.
In contrast, our \texttt{MDerive} abstraction is function-symbol agnostic, our three-tier oracle design supports recursive decomposition of compound keys, and our SA/QA attacker model anchors breakability to a concrete bit-level bound.

\noindent\textbf{Probabilistic model checking.}
Siedlecka-Lamch et al.~\cite{siedlecka2016probabilistic} developed a probabilistic model checker for security protocols.
Given a per-key compromise probability, the tool computes a risk tree and outputs a sensitivity assessment.
This line of work is complementary to ours: probabilistic tools quantify leakage risk but cannot produce the kind of concrete attack trace that a symbolic verifier produces, nor can they surface protocol-level multi-step attacks such as UMD.

\noindent\textbf{Summary.}
No prior work simultaneously supports protocol-agnostic modeling, compound-key decomposition, a graded bit-level attacker, and a non-intrusive implementation path; EntropyVerif is, to our knowledge, the first to check all four boxes, and is the first automated symbolic tool to discover a new protocol-level attack on Bluetooth Classic.
